\chapter{The Whole Machine, and Where It Goes Next}
\label{ch:wholemachine}

\chapterorient{This is the last chapter of the main text, and its job is to stand
back far enough to see everything at once. We have built, named, and composed every
part: the field theory, the calculus, the numerical engines, the six analyses, and
finally the synthesized library that renders them all as code. Here we tie the
single thread that runs through them---the \code{lifecycle} root from which every
other piece of the book is reachable---and draw the whole machine on one page. Then
we close two loops. The first is the book's opening promise (\cref{ch:bigidea}):
``fields as tensors, steps as pure functions'' was never an aesthetic preference; it
was the form in which the algorithm could be handed, unchanged, to an accelerator.
The second is the book's thesis itself, now that you have watched it proven component
by component: a large physics simulator is the repeated composition of a small set of
mathematical building blocks. We end honestly---naming the frontiers this book did
\emph{not} cross---and warmly, with what you can now do that you could not before.}

\section{One root, and everything below it}
\label{sec:wm-root}

Across five parts we disassembled a real electromagnetic simulator and, in Part~VI,
reassembled it as a small library: a handful of modules---\code{types},
\code{iteration}, \code{data-algebra}, \code{coordination}, and \code{drivers}---each
holding a few of the named machines we built (\cref{ch:synthlib},
\cref{ch:fivelibs}). \Cref{ch:composing} composed individual drivers from those
modules. One question remains, and it is the question that turns a \emph{library}
into a \emph{program}: what is the single entry point that ties all of it together?
What is the \code{main} of the synthesized library?

It has a name. It is \code{lifecycle}, the topmost composition in the \code{drivers}
module, and it is short enough to state in one line:
\begin{equation}
  \code{lifecycle} \;=\; \code{fold\_solve}\,\bigl(\code{dispatch}\,(\code{problem\_type}\;\code{cfg})\bigr)
                        \;\circ\; \code{build\_mesh}.
  \label{eq:wm-lifecycle}
\end{equation}
Read right to left: build the mesh from the parsed configuration; dispatch on the
configuration's \code{problem.type} to one of the six drivers; and wrap the chosen
driver in the adaptive estimate--mark--refine fold of \cref{ch:amr}, which threads
the mesh through repeated solves and---when adaptive refinement is switched
off---degenerates to a single solve. That is the entire top of the machine. Every
other piece of this book hangs beneath it.

\begin{keyidea}
The synthesized library has exactly one root: \code{lifecycle}. It is the \code{main}
of the whole simulator, and \cref{eq:wm-lifecycle} is its body in full---build the
mesh, dispatch on the problem type, fold the driver through the adaptive loop. From
this single function, \emph{every} component the book built is reachable by following
composition edges: the drivers, the solve corners, the assembly fold, the
matrix-free applies, the records. To hold \code{lifecycle} in your head is to hold
the whole machine.
\end{keyidea}

\subsection{The root, rendered as code}
\label{sec:wm-root-code}

\Cref{eq:wm-lifecycle} is the skeleton; here is the synthesized body, the way the
\code{drivers} module actually writes it. It is worth reading slowly, because every
identifier in it is a chapter.

\begin{lstlisting}[style=l4style]
-- The synthesized library's main. inputs = the parsed config; output = the
-- driver-selected physical product (capacitance | inductance | S-params |
-- (f,Q) | trajectory | fields | modes).
lifecycle :: IoData -> Product
lifecycle cfg =
  let mesh0 = build_mesh cfg                    -- (1) load + partition + a-priori refine   [ch:functionspaces]
      drv   = dispatch (problem_type cfg)       -- (2) select ONE driver: the specialization seam  [ch:lifecycle]
      step  = \m -> estimate_mark_refine cfg (drv cfg m)  -- one adaptive iterate: solve, then refine
  in  fold_solve step mesh0 (refinement cfg)    -- (3) adaptive estimate-mark-refine FOLD  [ch:amr]
  where
    -- the dispatch over the six driver defs, by problem type
    dispatch Electrostatic = electrostatic      -- [ch:electrostatics]
    dispatch Magnetostatic = magnetostatic      -- [ch:magnetostatics]
    dispatch Driven        = driven             -- [ch:driven]
    dispatch Transient     = transient          -- [ch:transient]
    dispatch Eigenmode     = eigenmode          -- [ch:eigenmodes]
    dispatch BoundaryMode  = boundary_mode      -- [ch:waveports]
\end{lstlisting}

\noindent Three lines of body and a six-way \code{dispatch}. Line~(1) is the mesh
build of \cref{ch:functionspaces}---the only stage every analysis runs identically.
Line~(2) is the \emph{specialization seam} of \cref{ch:lifecycle}: a single
field, \code{problem\_type}, read off the configuration record, selects which of the
six drivers fills in the one stage that differs. Line~(3) is the adaptive loop of
\cref{ch:amr}, written as exactly the \code{fold\_solve} combinator of \cref{ch:fold}:
its threaded value is the mesh, its body solves and then refines, and when no
refinement is requested it runs that body once. Nothing else is at the top. The
machine's entire control structure---which we spent Part~V tracing by hand---is these
three lines.

\begin{notationbox}
\code{dispatch} is an ordinary function from a \code{ProblemType} tag to a driver. It
is the synthesized form of the \texttt{switch} that \cref{ch:lifecycle} called the
specialization seam: above it, everything is shared (parse, mesh, the fold); below
it, the six drivers diverge. Because \code{dispatch} returns a \emph{function} (a
driver), \cref{eq:wm-lifecycle}'s middle factor is higher-order---the closure-returning
shape the calculus uses throughout (\cref{ch:language}). The whole of the
simulator's variety is funneled through this one tag.
\end{notationbox}

\subsection{Reachability: the root sees the whole book}
\label{sec:wm-reach}

The claim that ``everything is reachable from \code{lifecycle}'' is not rhetorical.
The synthesized library is a dependency graph: each function names the functions it
composes, and those edges form a finite tree rooted at \code{lifecycle}. Follow the
edges and you visit, in order, every machine the book built (\cref{fig:wm-root}).
\code{lifecycle} composes \code{build\_mesh} and \code{dispatch}; \code{dispatch}
reaches each of the six drivers; each driver composes an \emph{assemble} stage
(\code{fe\_assemble}, \cref{ch:assembly}), a \emph{solve} stage (one of four
coordination combinators, \cref{ch:situating}), and a \emph{reduce} stage (one of
three reduction verbs, \cref{ch:reductions}); the solve stages reach the Krylov
methods (\cref{ch:cg}, \cref{ch:gmres}), the multigrid preconditioner
(\cref{ch:multigrid}), and the smoothers (\cref{ch:smoothers}); and \emph{every}
operator application in any of them bottoms out in the same matrix-free apply
(\cref{ch:matrixfree}). There is no part of the book that this tree does not touch,
and no node in the tree that is not a chapter you have read.

\begin{figure}[p]
\centering
\begin{tikzpicture}[font=\footnotesize, >=Stealth, node distance=5mm and 4mm]
  % --- the root ---
  \node[stage, fill=simBgViolet, draw=simViolet, very thick, text width=30mm]
    (root) {\code{lifecycle}\\[1pt]\scriptsize the library's \code{main}};
  % parse / mesh feeding in
  \node[data, left=10mm of root, text width=16mm] (cfg)
    {\code{IoData}\\\scriptsize parsed config};
  \draw[flow] (cfg) -- node[above,font=\scriptsize]{parse} (root);
  % the AMR fold wrapping
  \node[opbox, above=7mm of root, text width=30mm] (amr)
    {\code{fold\_solve}: AMR loop\\\scriptsize estimate--mark--refine (\cref{ch:amr})};
  \draw[refedge] (amr.south) -- (root.north);
  \node[opbox, left=5mm of amr, text width=18mm] (mesh)
    {\code{build\_mesh}\\\scriptsize\cref{ch:functionspaces}};
  \draw[refedge] (mesh.south) to[out=-90,in=150] (root.north west);
  % --- dispatch fan-out to six drivers ---
  \node[stage, fill=simBgGray, draw=simGray, below=8mm of root, text width=24mm]
    (disp) {\code{dispatch} (\code{problem\_type})\\\scriptsize the specialization seam};
  \draw[flow] (root) -- (disp);
  % six driver boxes in a row
  \node[stage, below=8mm of disp, xshift=-44mm, text width=15mm] (d1)
    {\code{electro-}\\\code{static}};
  \node[stage, right=of d1, text width=15mm] (d2) {\code{magneto-}\\\code{static}};
  \node[stage, fill=simBgRust, draw=simRust, right=of d2, text width=13mm] (d3) {\code{driven}};
  \node[stage, right=of d3, text width=14mm] (d4) {\code{transient}};
  \node[stage, right=of d4, text width=15mm] (d5) {\code{eigen-}\\\code{mode}};
  \node[stage, right=of d5, text width=15mm] (d6) {\code{boundary-}\\\code{mode}};
  \foreach \d in {d1,d2,d3,d4,d5,d6}{\draw[flow] (disp) -- (\d);}
  % chapter tags under drivers
  \node[cornertag, below=0.5mm of d1] {\cref{ch:electrostatics}};
  \node[cornertag, below=0.5mm of d3] {\cref{ch:driven}};
  \node[cornertag, below=0.5mm of d5] {\cref{ch:eigenmodes}};
  % --- the shared body each driver composes ---
  \node[stage, fill=simBgBlue, draw=simBlue, below=10mm of d3, xshift=-14mm, text width=20mm]
    (asm) {\textbf{assemble}\\\code{fe\_assemble}\\\scriptsize\cref{ch:assembly}};
  \node[stage, fill=simBgRust, draw=simRust, right=7mm of asm, text width=24mm]
    (solve) {\textbf{solve}\\\scriptsize 4 corners (\cref{ch:situating})\\Krylov + multigrid};
  \node[stage, fill=simBgGold, draw=simGold, right=7mm of solve, text width=22mm]
    (red) {\textbf{reduce}\\\scriptsize 3 shapes (\cref{ch:reductions})};
  \foreach \d in {d1,d2,d3,d4,d5,d6}{
    \draw[refedge, simGray!60] (\d.south) to[out=-90,in=90] (asm.north);}
  \draw[flow] (asm) -- (solve);
  \draw[flow] (solve) -- (red);
  % the solve stack reaching down
  \node[opbox, below=6mm of solve, text width=26mm] (stack)
    {Krylov (\cref{ch:cg},\,\cref{ch:gmres})\\$\cdot$ V-cycle (\cref{ch:multigrid})\\$\cdot$ smoother (\cref{ch:smoothers})};
  \draw[refedge] (solve.south) -- (stack.north);
  \node[extern, below=5mm of stack, text width=30mm] (mf)
    {\textbf{matrix-free apply} (\cref{ch:matrixfree})\\\scriptsize every operator application bottoms out here};
  \draw[refedge] (stack.south) -- (mf.north);
  % records thread to the side
  \node[opbox, fill=white, right=8mm of stack, text width=20mm] (rec)
    {records: \code{SimState},\\\code{IoData}\\\scriptsize\cref{ch:records}, monad \cref{ch:monad}};
  \draw[refedge] (rec.west) -- (stack.east);
\end{tikzpicture}
\caption[The \code{lifecycle} root and its reachable machine]{The synthesized
library rooted at \code{lifecycle}, its \code{main}. The parsed \code{IoData} enters
on the left; \code{lifecycle} builds the mesh (\cref{ch:functionspaces}), wraps
everything in the adaptive \code{fold\_solve} loop (\cref{ch:amr}), and dispatches on
\code{problem\_type} (the specialization seam, \cref{ch:lifecycle}) to one of the six
drivers. Every driver composes the same three-stage body---\code{fe\_assemble}
(\cref{ch:assembly}), one of four solve corners (\cref{ch:situating}), one of three
reduction shapes (\cref{ch:reductions})---and every solve reaches the same nested
Krylov/multigrid/smoother stack, which bottoms out, like everything else, in the
matrix-free apply (\cref{ch:matrixfree}). The threaded state lives in the records
(\cref{ch:records}) carried by the solve monad (\cref{ch:monad}). Follow any edge and
you reach a chapter; this single tree \emph{is} the book.}
\label{fig:wm-root}
\end{figure}

\section{The whole machine on one page}
\label{sec:wm-onepage}

\Cref{fig:wm-root} drew the root and its first few layers. We now draw the
\emph{entire} machine---root to leaves, every driver, every combinator, every record,
every external-kernel boundary---in a single capstone figure
(\cref{fig:wm-hero}). It is dense by design: the point is not to read each box in
isolation but to \emph{see that everything connects}, that the simulator which looked,
from \cref{ch:targets}, like six elaborate and separate programs is in fact one finite
tree with a handful of shared leaves.

\begin{figure}[p]
\centering
\begin{tikzpicture}[font=\scriptsize, >=Stealth,
    grow=down, level distance=11mm, sibling distance=0mm,
    edge from parent/.style={draw=simGray, semithick, -{Stealth[length=1.4mm]}},
    every node/.style={align=center},
    level 1/.style={sibling distance=46mm},
    level 2/.style={sibling distance=23mm},
    level 3/.style={sibling distance=23mm}]
  \node[stage, fill=simBgViolet, draw=simViolet, very thick, text width=34mm]
    {\code{lifecycle}\\\scriptsize the whole machine (\S\ref{sec:wm-root})}
    child {node[stage, text width=22mm] {\code{build\_mesh}\\\scriptsize\cref{ch:functionspaces}}}
    child {node[opbox, text width=24mm] {\code{fold\_solve} (AMR)\\\scriptsize\cref{ch:amr}, \cref{ch:fold}}}
    child {node[stage, fill=simBgGray, draw=simGray, text width=26mm] {\code{dispatch}\\\scriptsize 6 drivers (\cref{ch:situating})}
      child {node[stage, text width=20mm] {\textbf{assemble}\\\code{fe\_assemble}\\\scriptsize\cref{ch:assembly}, \cref{ch:weak}}
        child {node[extern, text width=20mm] {matrix-free\\quadrature\\\scriptsize\cref{ch:matrixfree}}}
      }
      child {node[stage, fill=simBgRust, draw=simRust, text width=22mm] {\textbf{solve}\\\scriptsize 4 corners\\\code{solve\_family} /\\\code{frequency\_sweep} /\\\code{fold\_solve} / \code{eigsolve}}
        child {node[stage, fill=simBgBlue, draw=simBlue, text width=20mm] {Krylov + V-cycle\\\scriptsize\cref{ch:gmres}, \cref{ch:multigrid}}
          child {node[stage, fill=simBgGreen, draw=simGreen, text width=18mm] {smoother\\\scriptsize\cref{ch:smoothers}}
            child {node[extern, text width=18mm] {matrix-free\\apply\\\scriptsize\cref{ch:matrixfree}}}
          }
        }
        child {node[extern, text width=18mm] {\code{eigsolve}\\\scriptsize SLEPc kernel\\\cref{ch:krylovschur}}}
      }
      child {node[stage, fill=simBgGold, draw=simGold, text width=22mm] {\textbf{reduce}\\\scriptsize 3 shapes\\\code{gram\_reduce} /\\projection / table\\\cref{ch:reductions}}}
    };
  % the threaded-records side note
  \node[opbox, fill=white, text width=26mm, anchor=west] (rec) at (5.2,-1.1)
    {\textbf{records} (\cref{ch:records})\\\scriptsize \code{IoData} (read-only),\\\code{SimState} (threaded by\\the solve monad, \cref{ch:monad})};
\end{tikzpicture}
\caption[The whole machine on one page]{The capstone hero figure: the entire
synthesized simulator as a single finite tree, rooted at \code{lifecycle} and
expanded down to its leaves. Reading top to bottom: \code{lifecycle}
(\S\ref{sec:wm-root}) builds the mesh, wraps the run in the adaptive
\code{fold\_solve} loop (\cref{ch:amr}), and dispatches to one of six drivers
(\cref{ch:situating}). Each driver is the three-stage body
\emph{assemble--solve--reduce}: the assembly fold \code{fe\_assemble}
(\cref{ch:assembly}); one of the \emph{four} solve corners
(\code{solve\_family}, \code{frequency\_sweep}, \code{fold\_solve},
\code{eigsolve}), whose iterative cases nest a Krylov/multigrid/smoother stack
(\cref{ch:gmres}--\cref{ch:smoothers}) and whose black-box case calls the
\code{eigsolve} kernel (\cref{ch:krylovschur}); and one of the \emph{three} reduction
shapes (\cref{ch:reductions}). Every operator application in the iterative branch
bottoms out in the matrix-free apply (\cref{ch:matrixfree}, shown twice---it is the
universal leaf). The run's mutable state is the records (\cref{ch:records}) threaded
by the solve monad (\cref{ch:monad}). The figure is dense on purpose: it shows, in
one view, that the whole book is one connected machine.}
\label{fig:wm-hero}
\end{figure}

Two features of \cref{fig:wm-hero} deserve a word, because they are the whole point.
First, the tree is \emph{shallow and finite}. There is no recursion of structure, no
open-ended plumbing; from the root to the deepest leaf is six or seven composition
steps, and the same leaves are shared across all six drivers. A simulator is a small
tree, not a sprawl. Second, the leaves \emph{repeat}. The matrix-free apply appears
under assembly and again under every level of the solve stack; the reduction verbs are
shared across drivers that look physically unrelated; the records and the monad thread
through everything. The economy is not an accident of this book's exposition---it is a
property of the simulator, which is exactly why the book could be written at the length
it was rather than six times over.

\begin{keyidea}
The whole simulator fits on one page because it \emph{is} one finite, shallow tree of
shared parts. From \code{lifecycle} down to the matrix-free leaf is a handful of
composition steps, and the deepest, most-reused leaves---the matrix-free apply, the
reduction verbs, the records---are shared across all six drivers. The variety the
reader met in \cref{ch:targets} lives almost entirely in \emph{which few boxes a
driver swaps}, not in the structure of the tree.
\end{keyidea}

\section{Why this shape: the implementation view as a backend target}
\label{sec:wm-backend}

We can now close the book's oldest promise. Back in \cref{ch:bigidea} we made a
choice that, at the time, looked like a matter of taste: describe a simulation not as
in-place mutation of arrays but as pure functions over immutable tensors---``fields as
tensors, steps as pure functions.'' We gave three reasons (an explicit dependency
graph, clean composition, and a passing remark about accelerators) and promised the
third was the destination. This is that destination.

The synthesized library of \cref{fig:wm-hero} is written in a calculus whose semantics
were chosen to \emph{match} an external tensor backend
(\cref{ch:matrixfree}). Its values are immutable tensors. Its operations are pure
functions: an operator apply takes an input tensor and returns an output tensor, with
no aliasing, no hidden global, and no ordering hazard to police. Its combinators---the
fold, the map, the solve corners---express coordination without ever naming a memory
cell. This is precisely the form a GPU, or any massively parallel backend, wants to
consume: ``here is an immutable input tensor; produce an output tensor,'' applied to
all entries at once (\cref{fig:wm-backend}). The matrix-free apply at the bottom of
the tree is a sequence of element-local tensor contractions
(\cref{ch:matrixfree})---a pure, whole-tensor transformation that lowers onto an
accelerator kernel without rewriting.

\begin{figure}[t]
\centering
\begin{tikzpicture}[font=\footnotesize, >=Stealth, node distance=8mm and 12mm]
  % left: the synthesized library (pure tensor code)
  \node[stage, fill=simBgViolet, draw=simViolet, text width=34mm, align=left] (lib)
    {\textbf{synthesized library}\\[2pt]\scriptsize
     \lstinline[style=l4style]|apply :: Tensor[N] -> Tensor[N]|\\
     \scriptsize immutable tensors,\\pure functions, combinators\\(\cref{ch:bigidea}, \cref{fig:wm-hero})};
  % middle: the lowering arrow
  \node[opbox, right=of lib, text width=22mm] (lower)
    {\textbf{lower}\\\scriptsize no rewrite:\\the form already matches};
  % right: the accelerator
  \node[extern, right=of lower, text width=30mm, minimum height=18mm] (gpu)
    {\textbf{GPU / tensor backend}\\[2pt]\scriptsize whole-tensor ops,\\no aliasing,\\all entries at once\\(\cref{ch:matrixfree})};
  \draw[flow] (lib) -- (lower);
  \draw[flow] (lower) -- (gpu);
  % the closed-loop callback to ch:bigidea
  \node[font=\scriptsize\itshape, text=simGray, below=3mm of lower, text width=44mm, align=center]
    {the promise of \cref{ch:bigidea}, kept:\\``fields as tensors, steps as pure functions''\\was always heading here};
  % grid of cells on the GPU to suggest parallelism
  \begin{scope}[shift={($(gpu.south)+(0,-1.0)$)}]
    \foreach \i in {0,1,2,3,4,5}{
      \fill[simBgViolet, draw=simViolet] (\i*0.42-1.25,0) rectangle (\i*0.42-0.9,0.32);}
    \node[font=\scriptsize, text=simViolet] at (0,-0.32) {one kernel, all cells at once};
  \end{scope}
\end{tikzpicture}
\caption[The implementation view as a backend-lowering target]{Why the pure picture
was the right one. The synthesized library (left) is written in tensor-in,
tensor-out pure functions---\lstinline[style=l4style]|apply :: Tensor[N] -> Tensor[N]|,
no mutation, no aliasing. That is, almost word for word, the form a GPU or other
tensor backend consumes (right): a whole-tensor operation applied to every entry at
once. So lowering the library onto an accelerator is not a translation---the
description \emph{already matches} the hardware's model. This closes the promise of
\cref{ch:bigidea}: ``fields as tensors, steps as pure functions'' was never an
aesthetic choice; it was the form in which the algorithm could be handed to an
accelerator and run unchanged.}
\label{fig:wm-backend}
\end{figure}

This is the deep reason the implementation view of Part~VI is worth having as a
distinct rendering at all. It is not merely ``the same spec, written as code.'' It is
the spec written in \emph{the form that runs on the hardware the whole effort
targets}. The book never produces a port; it produces this layered description. But the
top layer of that description---the synthesized library---is deliberately shaped so
that the final step, lowering it onto burn's device backends or any equivalent
tensor library, is the short one. The hard work of saying \emph{what the algorithm
is}, free of any one machine's cost model, is what the five parts before this one did.

\begin{keyidea}
The functional rendering matters because it is the \emph{backend-lowering target}. The
synthesized library's semantics---immutable tensors, pure functions, whole-tensor
combinators---were chosen to match an external GPU/tensor backend, so the algorithm can
be handed to an accelerator and run \emph{unchanged}. ``Fields as tensors, steps as
pure functions'' (\cref{ch:bigidea}) was never decoration; it was a commitment to the
one description from which the accelerator form is a short step rather than a rewrite.
\end{keyidea}

\subsection{A worked dispatch: the final \emph{everything-connects} trace}
\label{sec:wm-worked}

We have one worked example left, and it is the whole book in miniature: take the
\code{lifecycle} root, feed it one real configuration, and trace the single dispatch
all the way down until it bottoms out in the machinery of every part.

\begin{workedexample}{One dispatch, root to leaf}{wm-dispatch}
An engineer hands the simulator a configuration file for a capacitance extraction: a
mesh of two conductors, the permittivity $\varepsilon$ of the dielectric between them,
the terminals to drive, and the field \code{problem\_type:~Electrostatic}. We trace
\code{lifecycle cfg} from the top.

\medskip
\textbf{(1) Build the mesh.} \code{lifecycle} calls \code{build\_mesh cfg}
(\cref{ch:functionspaces}), which loads the geometry, partitions it, and returns the
\code{Mesh}---the first value the run \emph{produces} rather than reads. This stage is
identical for all six analyses.

\textbf{(2) Dispatch.} \code{lifecycle} reads \code{problem\_type cfg}, finds
\code{Electrostatic}, and \code{dispatch} returns the \code{electrostatic} driver
(\cref{ch:electrostatics}). One field, one branch, the entire choice of \emph{which}
analysis runs. Had the field read \code{Driven}, the same line would have returned the
\code{driven} driver of \cref{wex:trace-driven}; the root does not change.

\textbf{(3) The adaptive fold.} \code{lifecycle} wraps the chosen driver in
\code{fold\_solve} (\cref{ch:amr}, \cref{ch:fold}). The configuration requests no
refinement, so the fold runs its body \emph{once}---the straight-line lifecycle of
\cref{ch:lifecycle}. Everything below is one pass.

\textbf{(4) Inside the driver: assemble.} \code{electrostatic cfg} lays the nodal
$\Hone$ space (\cref{ch:functionspaces}) and assembles the stiffness operator $\Kop$
by the weak-form fold \code{fe\_assemble} (\cref{ch:assembly}, \cref{ch:weak})---one
term, the $\varepsilon$-weighted gradient---with the essential boundary degrees of
freedom eliminated (\cref{ch:bc}). $\Kop$ is symmetric positive-definite and built
\emph{once}.

\textbf{(5) Inside the driver: solve.} The driver maps the fixed-operator solve over
the per-terminal right-hand sides: \code{solve\_family} (corner~1 of
\cref{ch:situating}). Each solve is CG (\cref{ch:cg}, because $\Kop$ is SPD)
preconditioned by a multigrid V-cycle (\cref{ch:multigrid}) over a smoother
(\cref{ch:smoothers}), and \emph{every} operator application inside all of them is the
matrix-free apply (\cref{ch:matrixfree})---the same element-local tensor contraction
that the driven filter's GMRES bottomed out in. The solver's running state lives in the
\code{SimState} record (\cref{ch:records}) threaded by the solve monad
(\cref{ch:monad}).

\textbf{(6) Inside the driver: reduce.} The driver reduces the solved family by the
symmetric Gram \code{gram\_reduce} (\cref{ch:reductions}), $C_{ij}=V_j^{\!\top}\Kop
V_i$, pairing the fields with themselves through the energy operator to produce the
$n\times n$ capacitance matrix.

\textbf{(7) Output.} \code{lifecycle} writes the product and the run metadata
(\cref{ch:lifecycle}).

\medskip
Read the chapter tags down the trace: \cref{ch:functionspaces}, \cref{ch:lifecycle},
\cref{ch:amr}, \cref{ch:assembly}, \cref{ch:bc}, \cref{ch:situating}, \cref{ch:cg},
\cref{ch:multigrid}, \cref{ch:smoothers}, \cref{ch:matrixfree}, \cref{ch:records},
\cref{ch:monad}, \cref{ch:reductions}. \emph{One dispatch touched every part of the
book.} And the only lines that distinguish this run from the driven filter of
\cref{wex:trace-driven} are the three swapped cells of \cref{ch:fulltrace}---the
function space, the solve corner, and the reduction. The root is the same; the path
down from it is the whole machine; and which machine you get is a single field in a
file.
\end{workedexample}

\section{Where it goes next: the deferred frontiers}
\label{sec:wm-frontiers}

A closing chapter owes the reader an honest accounting of what the book did
\emph{not} do. There are three frontiers the synthesized library stops at, and it is
worth being precise about each---not as failures, but as the well-marked edges of a
deliberately scoped effort, each with a known direction onward (\cref{fig:wm-next}).

\paragraph{Distributed execution and sharding.}
The entire book targets a \emph{single machine}: one CPU lifting onto one GPU through a
device backend (\cref{ch:targets}). Real Palace runs across many machines, partitioning
the mesh and the fields into shards and exchanging data between them by message passing.
We deliberately read every distributed type as its single-rank equivalent and set the
parallel machinery aside, because the sharding theory assumes a memory and lifetime
structure that the pure, immutable-tensor spine of this book has rewritten---lifting it
naively could destabilize the very abstractions that made the spec clean. The
\emph{mathematics} of decomposition---splitting a tensor-field problem into coupled
sub-problems on sub-domains---is a genuine and tractable next direction, and it composes
with the pure picture rather than fighting it: a sharded solve is the same combinators
over a tensor that happens to be distributed. But it is future work, gated behind a
careful check that it does not disturb the single-machine spine, and we have marked it
as such throughout.

\paragraph{The opaque kernels, fully realized.}
Three places in the synthesized library are \code{\#extern} boundaries: a named
type signature standing in for an implementation we describe but do not fully render.
Each is an opaque library kernel that the spine \emph{depends on}, and each has a
known implementation in the vocabulary we built. The libCEED element-quadrature kernel
inside \code{fe\_assemble} is matrix-free operator application as element-local tensor
contractions (\cref{ch:matrixfree})---the implementation is understood; only its full
rendering is deferred. The SLEPc eigenvalue loop behind \code{eigsolve} is a
constructive Krylov--Schur iteration in our own \code{krylov\_step} and
\code{lanczos\_step} vocabulary (\cref{ch:krylovschur}). The MFEM time-step behind the
transient \code{fold\_solve} is one ODE-integrator step (\cref{ch:time}). For each, the
book gives \emph{both} surfaces---the opaque contract the spine calls, and the
in-our-vocabulary implementation it corresponds to---and leaves the final reconciliation
of the two as marked, well-understood work.

\paragraph{Deeper modularization.}
The five-module partition of the synthesized library (\cref{ch:fivelibs}) is a sensible
default, not a theorem. As the library is used---as components are reused across
drivers, as new analyses are added---some boundaries will want to move: a type may
migrate to cluster with the API that consumes it; a combinator family may earn its own
module. This is the ordinary refinement-by-use of any growing library, and the book
leaves the partition explicitly open to it.

\begin{figure}[t]
\centering
\begin{tikzpicture}[font=\footnotesize, >=Stealth, node distance=6mm and 10mm]
  % the firm core
  \node[stage, fill=simBgViolet, draw=simViolet, very thick, text width=30mm, minimum height=16mm]
    (core) {\textbf{this book}\\[1pt]\scriptsize the single-machine\\synthesized library\\(\cref{fig:wm-hero})};
  % three frontiers radiating out
  \node[extern, above right=4mm and 16mm of core, text width=34mm] (shard)
    {\textbf{distributed / sharding}\\\scriptsize decomposition as math;\\gated on spine stability\\(\cref{ch:targets})};
  \node[extern, right=22mm of core, text width=34mm] (kern)
    {\textbf{opaque kernels, realized}\\\scriptsize libCEED quadrature (\cref{ch:matrixfree}),\\\code{eigsolve} (\cref{ch:krylovschur}),\\ODE step (\cref{ch:time})};
  \node[extern, below right=4mm and 16mm of core, text width=34mm] (mod)
    {\textbf{deeper modularization}\\\scriptsize refine the 5-module\\partition by use (\cref{ch:fivelibs})};
  \draw[refedge] (core) -- (shard);
  \draw[refedge] (core) -- (kern);
  \draw[refedge] (core) -- (mod);
  \node[font=\scriptsize\itshape, text=simGray, below=2mm of core] {where it goes next};
\end{tikzpicture}
\caption[The deferred frontiers]{Where the book goes next. The firm single-machine
synthesized library (left) is bordered by three marked frontiers, each a direction
onward rather than a gap. \emph{Distributed/sharding}: the mathematics of domain
decomposition, composing with the pure picture, gated behind a check that it does not
destabilize the single-machine spine. \emph{The opaque kernels realized}: the three
\code{\#extern} boundaries---libCEED quadrature (\cref{ch:matrixfree}), the
\code{eigsolve} loop (\cref{ch:krylovschur}), the ODE step (\cref{ch:time})---each with
a known in-our-vocabulary implementation to be fully rendered. \emph{Deeper
modularization}: the five-module partition refined by use. None is a failure; each is a
well-understood next step from a deliberately scoped foundation.}
\label{fig:wm-next}
\end{figure}

\begin{pitfall}
It would be a misreading to take the \code{\#extern} boundaries as holes in the
argument. They are the opposite: a \code{\#extern} is a place where the book knows
\emph{exactly} what the implementation is---it has a type, a contract, and a
described realization in the calculus---and has chosen to defer the full rendering, not
the understanding. The frontier is one of \emph{effort remaining}, not of
\emph{comprehension missing}. That distinction is the difference between an honest
boundary and a gap.
\end{pitfall}

\section{The reflective close: three views, one machine}
\label{sec:wm-reflect}

Step all the way back. This book showed you the \emph{same} simulator three times.

Part~V showed it \emph{top-down}: the six analyses as the engineer meets them, each a
physical question with a physical product, walked through the lifecycle stage by stage
until the whole thing collapsed into one composition (\cref{ch:fulltrace}). Parts~II--IV
showed it \emph{bottom-up}: the vocabulary, built from the ground---the tensor and its
shapes, the de Rham complex, the weak forms, the calculus of pure transformations, the
Krylov and multigrid engines---each piece firm in itself before it was composed.
Part~VI showed it as the \emph{synthesized library}: the same machines, written out
with their types and bodies, rendered as the code an accelerator would run. Three
views---top-down modalities, bottom-up vocabulary, the implementation
library---and the central claim of the book is that they are \emph{one machine seen
three ways} (\cref{fig:wm-triality}).

\begin{figure}[t]
\centering
\begin{tikzpicture}[font=\footnotesize, >=Stealth]
  % the central machine
  \node[stage, fill=simBgViolet, draw=simViolet, very thick, circle, minimum size=22mm,
        text width=18mm, align=center] (m) at (0,0) {\textbf{one}\\\textbf{machine}};
  % three views around it
  \node[stage, fill=simBgGold, draw=simGold, text width=30mm, align=center] (top) at (0,2.7)
    {\textbf{top-down: modalities}\\\scriptsize 6 analyses, walked\\(Part~V, \cref{ch:fulltrace})};
  \node[stage, fill=simBgBlue, draw=simBlue, text width=30mm, align=center] (bot) at (-4.4,-2.2)
    {\textbf{bottom-up: vocabulary}\\\scriptsize tensors, de Rham, calculus,\\engines (Parts~II--IV)};
  \node[stage, fill=simBgGreen, draw=simGreen, text width=30mm, align=center] (lib) at (4.4,-2.2)
    {\textbf{the library: code}\\\scriptsize synthesized, accelerator-ready\\(Part~VI)};
  \draw[depedge] (top) -- (m);
  \draw[depedge] (bot) -- (m);
  \draw[depedge] (lib) -- (m);
  % the "they agree" annotations on the connecting edges
  \node[font=\scriptsize\itshape, text=simGray] at (1.9,1.5) {same composition};
  \node[font=\scriptsize\itshape, text=simGray] at (-2.9,-0.6) {same parts};
  \node[font=\scriptsize\itshape, text=simGray] at (2.9,-0.6) {same algorithm};
\end{tikzpicture}
\caption[The three views as one machine]{The triality. The book renders one
simulator three ways: \emph{top-down} as the six analyses walked through the lifecycle
(Part~V); \emph{bottom-up} as the vocabulary built from tensors and the de Rham complex
up through the calculus and the numerical engines (Parts~II--IV); and as the
\emph{synthesized library}, the same machines written as accelerator-ready code
(Part~VI). The book's central claim is that the three are not three descriptions of
three things but three views of \emph{one machine}---they compose the same parts into the
same algorithm, and \cref{ch:fulltrace} caught all three agreeing in a single trace.}
\label{fig:wm-triality}
\end{figure}

And here is the deep claim those three views were built to establish, now that you have
seen it proven component by component rather than asserted: \emph{a large physics
simulator is the repeated composition of a small set of mathematical building blocks.}
Not a sprawl of special cases. A small kit, composed (\cref{fig:wm-kit}). The kit is
short enough to list:

\begin{itemize}[leftmargin=*]
  \item \textbf{Four solve corners} (\cref{ch:situating})---the fixed-operator map, the
  operator-varying map, the state fold, and the black-box eigensolve. Every solve in
  every analysis is one of these four ways of threading the step.
  \item \textbf{Three reduction shapes} (\cref{ch:reductions})---the symmetric Gram, the
  port-projection, and the per-mode table. Every physical product is one of these three
  ways of collapsing a solution family to numbers.
  \item \textbf{The de Rham complex} (\cref{ch:derham})---the four function spaces
  ($\Hone$, $\Hcurl$, $\Hdiv$, $\Ltwo$) and the gradient, curl, and divergence that link
  them, from which every operator in the book is assembled.
  \item \textbf{The shared substrate}---the fold (\cref{ch:fold}), the state monad
  (\cref{ch:monad}), and the matrix-free operator (\cref{ch:matrixfree}). These thread
  through everything: the fold coordinates, the monad owns the state, the matrix-free
  apply is the universal leaf.
\end{itemize}

\begin{figure}[t]
\centering
\begin{tikzpicture}[font=\footnotesize, >=Stealth, node distance=4mm and 6mm]
  % four corners (2x2)
  \node[font=\small\bfseries, text=simRust] at (-3.2,2.4) {4 solve corners};
  \node[stage, fill=simBgRust, draw=simRust, minimum width=16mm, minimum height=8mm] (c1) at (-4.0,1.5) {\scriptsize fixed map};
  \node[stage, fill=simBgRust, draw=simRust, minimum width=16mm, minimum height=8mm] (c2) at (-2.3,1.5) {\scriptsize op-vary map};
  \node[stage, fill=simBgRust, draw=simRust, minimum width=16mm, minimum height=8mm] (c3) at (-4.0,0.6) {\scriptsize state fold};
  \node[extern, minimum width=16mm, minimum height=8mm] (c4) at (-2.3,0.6) {\scriptsize black-box};
  \node[cornertag] at (-3.15,0.0) {\cref{ch:situating}};
  % three reductions
  \node[font=\small\bfseries, text=simGold] at (0.9,2.4) {3 reductions};
  \node[stage, fill=simBgGold, draw=simGold, minimum width=20mm, minimum height=6mm] (r1) at (0.9,1.6) {\scriptsize Gram};
  \node[stage, fill=simBgGold, draw=simGold, minimum width=20mm, minimum height=6mm] (r2) at (0.9,0.9) {\scriptsize projection};
  \node[stage, fill=simBgGold, draw=simGold, minimum width=20mm, minimum height=6mm] (r3) at (0.9,0.2) {\scriptsize per-mode table};
  \node[cornertag] at (0.9,-0.35) {\cref{ch:reductions}};
  % shared substrate
  \node[font=\small\bfseries, text=simBlue] at (4.4,2.4) {shared substrate};
  \node[opbox, minimum width=24mm] (s1) at (4.4,1.6) {\scriptsize the fold (\cref{ch:fold})};
  \node[opbox, minimum width=24mm] (s2) at (4.4,0.9) {\scriptsize the monad (\cref{ch:monad})};
  \node[extern, minimum width=24mm] (s3) at (4.4,0.2) {\scriptsize matrix-free (\cref{ch:matrixfree})};
  % de Rham underneath
  \node[space, minimum width=86mm, minimum height=8mm] (dr) at (0.5,-1.3)
    {\footnotesize \textbf{the de Rham complex}: $\Hone \xrightarrow{\Grad} \Hcurl \xrightarrow{\Curl} \Hdiv \xrightarrow{\Div} \Ltwo$ \quad(\cref{ch:derham})};
  % composition arrows down to a product
  \draw[flow] (c4.south) -- (dr.north);
  \draw[flow] (r3.south) -- (dr.north);
  \draw[flow] (s3.south) -- (dr.north);
\end{tikzpicture}
\caption[The small kit of building blocks]{The book's thesis in one figure: a large
simulator is the repeated composition of a small kit. \emph{Four solve corners}
(\cref{ch:situating}) $\times$ \emph{three reduction shapes} (\cref{ch:reductions}),
resting on a \emph{shared substrate}---the fold (\cref{ch:fold}), the state monad
(\cref{ch:monad}), and the matrix-free operator (\cref{ch:matrixfree})---all assembled
over the \emph{de Rham complex} (\cref{ch:derham}). Six physically distinct analyses are
choices of one corner and one reduction over this same kit. The kit is small; the
simulator is its composition.}
\label{fig:wm-kit}
\end{figure}

That is the entire kit. Four corners, three reductions, one de Rham complex, and a
shared substrate of fold, monad, and matrix-free operator. Six physically distinct
simulations---capacitance, inductance, resonances, S-parameters, time-domain fields,
waveguide modes---are each a choice of \emph{one corner and one reduction} over this
kit, assembled on \emph{one} function-space complex. The thesis we held only as a
promise in \cref{ch:targets}---``six settings of one machine''---is now a fact you have
watched assemble itself, block by block, across forty-eight chapters.

\begin{keyidea}
A large physics simulator is the repeated composition of a small set of mathematical
building blocks---four solve corners, three reduction shapes, the de Rham complex, the
fold, the monad, and the matrix-free operator. You have now seen the whole machine: one
root, one finite tree, a handful of shared leaves. And you have seen \emph{why} it is
shaped the way it is---pure functions over immutable tensors---because that is the form
in which it will run on the accelerators it was always meant for. The variety is in the
composition; the kit is small.
\end{keyidea}

\subsection{What you can now do}
\label{sec:wm-cando}

The point of building this understanding was never the disassembly of one simulator. It
was a way of reading. With the kit in hand, you can do three things you could not before.

You can \emph{read any simulator's algorithms this way}. Open an unfamiliar
finite-element or finite-difference code and the questions are no longer ``what is all
this?'' but the three of \cref{ch:bigidea}: what operations happen, who owns the state,
how is the evolution coordinated? The in-place loops resolve into pure transformations;
the tangle of arrays resolves into threaded values; the control flow resolves into folds
and maps.

You can \emph{recognize the corners}. Faced with a new analysis, you can ask which of
the four solve corners its middle stage occupies---is the operator fixed or varying, is
the solve a map or a fold or a black box?---and which of the three reduction shapes
collapses its answer. Most of a new simulator is already familiar before you read a line
of its solver, because most of it is the shared substrate.

And you can \emph{compose a new analysis}. To add an analysis the book did not
cover---a thermal solve, a coupled multiphysics problem---is to choose a function space
from the de Rham complex, assemble its operators by the weak-form fold, pick the solve
corner its physics demands, and reduce its solution to the product an engineer wants. The
machine is a kit; you now know the pieces and how they snap together.

\section{A pointer to the appendices}
\label{sec:wm-appendix}

The figures of this book are deliberately pictures---block diagrams, composition trees,
trace tables---because pictures are how the \emph{shape} of the machine is best seen. A
reader who wants the exact rules behind the pictures will find them in the appendices,
which carry the deep computation semantics the main text kept at arm's length: the formal
calculus and its grammar (the precise syntax of the language sketched in
\cref{ch:language}), the shape algebra and its named-axis congruence rules, the algebraic
laws each combinator obeys, the numerical conventions that distinguish a transparent
performance trick from a load-bearing one, and the full matrix-free kernel detail behind
the universal leaf of \cref{ch:matrixfree}. The main text is the machine drawn; the
appendices are the machine specified to the last rule. Read them when a picture leaves you
wanting the proof.

\begin{keyidea}
\textbf{You have now seen the whole machine}---and why it is shaped for the accelerators it
will run on. A large simulator is not a monolith but the repeated composition of a small
kit of mathematical building blocks, rooted in a single \code{lifecycle} function from which
every part is reachable. The pure, immutable-tensor form was the right choice all along,
because it is the form in which the algorithm becomes a backend-lowering target rather than
a rewrite. That is the end of the disassembly, and the beginning of a way of reading every
simulator you meet from here on.
\end{keyidea}

\medskip
\noindent And that is the whole machine. We met it as six black boxes in
\cref{ch:targets}; we leave it as one small, connected, accelerator-ready
composition---understood, on one page, root to leaf. Thank you for building it with us.

\summarysection
This closing chapter views the whole simulator at once. The synthesized library has a
single root, \code{lifecycle}---its \code{main}---whose body is
\code{lifecycle $=$ fold\_solve (dispatch (problem\_type cfg)) $\circ$ build\_mesh}: build
the mesh (\cref{ch:functionspaces}), dispatch on the problem type at the specialization
seam (\cref{ch:lifecycle}) to one of six drivers, and fold the chosen driver through the
adaptive estimate--mark--refine loop (\cref{ch:amr}, degenerating to a single solve when
refinement is off). Every component of the book is reachable from this root by following
composition edges, and the whole machine fits on one page as a finite, shallow tree whose
deepest, most-reused leaf is the matrix-free apply (\cref{ch:matrixfree}). The functional
rendering is the \emph{backend-lowering target}: its immutable-tensor, pure-function,
whole-tensor-combinator form matches an external GPU/tensor backend, so the algorithm
lowers onto an accelerator unchanged---closing the promise of \cref{ch:bigidea}. Three
frontiers remain, marked as directions onward rather than failures: distributed/sharding
(deferred, gated on spine stability); the three \code{\#extern} opaque kernels (libCEED
quadrature, the \code{eigsolve} loop, the ODE step---each with a known in-vocabulary
implementation); and deeper modularization. Reflectively, the book showed one simulator
three ways---top-down modalities (Part~V), bottom-up vocabulary (Parts~II--IV), and the
synthesized library (Part~VI)---and proved they are one machine. The deep thesis, now
demonstrated component by component: a large physics simulator is the repeated composition
of a small kit---four solve corners (\cref{ch:situating}), three reduction shapes
(\cref{ch:reductions}), the de Rham complex (\cref{ch:derham}), and the shared substrate of
the fold (\cref{ch:fold}), the monad (\cref{ch:monad}), and the matrix-free operator
(\cref{ch:matrixfree}). With that kit, the reader can read any simulator's algorithms this
way, recognize its solve corners, and compose a new analysis. The appendices carry the exact
computation semantics behind the pictures.

\furtherreading
The simulator dissected throughout this book is AWS Labs' Palace~\citep{palace2023}, whose
source is the ground truth every layer of the spec cites. The matrix-free, element-local
operator-evaluation model that the synthesized library lowers onto an accelerator---the
universal leaf of \cref{fig:wm-hero} and the libCEED \code{\#extern} of
\S\ref{sec:wm-frontiers}---is developed by the CEED project~\citep{ceed2021}, the standard
reference for high-order matrix-free finite-element kernels. The iterative-solver
machinery the solve corners nest---the Krylov methods, the multigrid preconditioners, and
the reuse economics across an operator family---is treated comprehensively by
Saad~\citep{saad2003}, on which \cref{ch:gmres}, \cref{ch:multigrid}, and
\cref{ch:krylovschur} all draw. For the frontier of distributed execution and domain
decomposition (\S\ref{sec:wm-frontiers}), Saad's later chapters and the parallel-solver
literature are the natural next reading; the mathematics composes with the pure picture
this book built, and is the clearest single direction onward.

\exercisessection
\begin{enumerate}[nosep]
  \item \textbf{Read the root.} Write out the body of \code{lifecycle} from memory
  (\cref{eq:wm-lifecycle}), then name the chapter responsible for each of its three stages
  and for the six-way \code{dispatch}. Which single stage is identical for all six
  analyses, and which single field selects among them?
  \item \textbf{Trace a different dispatch.} Repeat the worked trace of
  \cref{wex:wm-dispatch}, but for a configuration with \code{problem\_type:~Eigenmode}.
  Which line of \code{lifecycle} changes its result, which driver does \code{dispatch}
  return, and at what point does this trace \emph{differ} from the electrostatic one---the
  solve corner, the reduction, or both? (Cross-check against \cref{ch:eigenmodes}.)
  \item \textbf{The universal leaf.} \Cref{fig:wm-hero} shows the matrix-free apply
  appearing in more than one place in the tree. List every distinct path from the root
  \code{lifecycle} down to a matrix-free apply, and explain why \emph{every} operator
  application in an iterative solve---the Krylov step, the smoother sweep, the residual---
  bottoms out in the same leaf.
  \item \textbf{Why pure, again.} In one paragraph, explain to a skeptic who has written
  fast in-place C why the synthesized library is written in pure functions over immutable
  tensors rather than mutating arrays. Your answer should connect the choice to the
  accelerator-lowering claim of \S\ref{sec:wm-backend} and the dependency-graph and
  composition payoffs of \cref{ch:bigidea}---and should address the objection that purity
  ``copies the array on every step.'' (Hint: revisit the pitfall of \cref{ch:bigidea}.)
  \item \textbf{The kit, counted.} \Cref{fig:wm-kit} lists the kit: four solve corners,
  three reduction shapes, the de Rham complex, and a shared substrate. For each of the six
  analyses of \cref{ch:targets}, state which solve corner and which reduction shape it uses.
  How many of the twelve corner-$\times$-reduction combinations does Palace actually
  occupy, and what does the answer say about the economy of the design?
  \item \textbf{Where would sharding go?} (Open-ended.) Suppose you were to add distributed
  execution to the synthesized library (\S\ref{sec:wm-frontiers}). Identify the \emph{one}
  node in the tree of \cref{fig:wm-hero} where a domain decomposition would most naturally
  enter, and argue whether it would change the \emph{structure} of the tree or only the
  implementation of a leaf. Why does the book claim sharding ``composes with the pure
  picture rather than fighting it''? What is the risk that gates it as future work?
  \item \textbf{Compose a new analysis.} (Open-ended.) Sketch how you would add a
  steady-state \emph{thermal} analysis (solve $-\Div(\kappa\Grad T)=q$ for a temperature
  field $T$, reduce to a thermal-resistance matrix) to the synthesized library. Which
  function space of the de Rham complex (\cref{ch:derham}) does $T$ live in? Which existing
  solve corner (\cref{ch:situating}) does the steady linear solve occupy? Which reduction
  shape (\cref{ch:reductions}) produces the resistance matrix, and which existing driver is
  your new one most nearly a copy of? Which pieces of the kit would you reuse unchanged, and
  what is the smallest set of new pieces you would have to write?
\end{enumerate}
